%% LyX 2.4.4 created this file.  For more info, see https://www.lyx.org/.
%% Do not edit unless you really know what you are doing.
\documentclass[english]{article}
\usepackage[T1]{fontenc}
\usepackage[latin9]{inputenc}
\usepackage{float}
\usepackage{units}
\usepackage{enumitem}
\usepackage{amsmath}
\usepackage{amsthm}
\usepackage{graphicx}
\usepackage{geometry}
\geometry{verbose,tmargin=1in,bmargin=1in,lmargin=1in,rmargin=1in}

\makeatletter
%%%%%%%%%%%%%%%%%%%%%%%%%%%%%% Textclass specific LaTeX commands.
\numberwithin{equation}{section}
\numberwithin{figure}{section}
\newlength{\lyxlabelwidth}      % auxiliary length 

%%%%%%%%%%%%%%%%%%%%%%%%%%%%%% User specified LaTeX commands.
\usepackage{fancyhdr}
\usepackage{lastpage}
\pagestyle{fancy}
\usepackage{enumitem}
\usepackage{ifthen}
%sets math fonts to be more readable
\everymath=\expandafter{\the\everymath\displaystyle}
%\DeclareMathSizes{10}{10}{10}{10}
\usepackage{multicol}
% Added by lyx2lyx
\setlength{\parskip}{\smallskipamount}
\setlength{\parindent}{0pt}

\makeatother

\usepackage{babel}
\begin{document}
\renewcommand{\author}{Dr.\,James Ashe}

\newcommand{\classNum}{331}

\newcommand{\classSec}{A}

\newcommand{\nameType}{Name:}

\newcommand{\examType}{Take home quiz}

\renewcommand{\date}{Due November 11, 2012}

%%First page's header%%%%%%%%%%%%%%%%%%%%%%
\fancypagestyle{firststyle} %%header settings for first pagr
{
	\fancyhead[L]{MTH \classNum{}(\classSec{}) \author{}\\
	\textbf{\examType{}} \date{}}
	\fancyhead[C]{\textbf{\nameType}}
	\setlength{\headsep}{14pt}
}
%%%%%%%%%%%%%%%%%%%%%%%%%%%%%%%%%%%%%%%%%%

%%Next page's header%%%%%%%%%%%%%%%%%%%%%%%
\setlist{nolistsep}
\lhead{\classNum{}(\classSec{})} 
\chead{\textbf{\nameType}} 
\cfoot{\thepage{}~of~\pageref{LastPage}} 
\setlength{\headsep}{5pt} 
%%%%%%%%%%%%%%%%%%%%%%%%%%%%%%%%%%%%%%%%%%%

%%Various settings; see the preamble for other settings%%%%%
%%\setenumerate[0]{leftmargin=12pt,itemindent=0pt,labelwidth=0pt}
\renewcommand{\labelenumi}{\textbf{\arabic{enumi}.}}
\renewcommand{\labelenumii}{\textbf{\alph{enumii}.}}
\thispagestyle{firststyle}
%%%%%%%%%%%%%%%%%%%%%%%%%%%%%%%%%%%%%%%%%%

\section*{10.1: Approximations with Taylor Polynomials. }

\textbf{HW problems 21-22 and 33-36.}\emph{ Approximate the given
quantities using Taylor polynomials with $n=3$, and compute the absolute
error in the approximation assuming the exact value is given by a
calculator.}
\begin{enumerate}
\item $e^{0.2}$\vfill
\item $\ln(1.01)$\vfill
\item $\sin(\frac{\pi+.02}{2})$ use center $\frac{\pi}{2}$.\vfill
\end{enumerate}

\section*{10.2: Find the Interval of Convergence. }

\textbf{HW problems 9-10 and 21-24}\textbf{\emph{.}}\emph{ Determine
the radius of convergence of the following power series. The test
the endpoints to determine the interval of convergence. }
\begin{enumerate}[resume]
\item ${\displaystyle \sum\left(\frac{3x}{2}\right)^{k}}$\vfill
\item $\sum\frac{k^{2}x^{2k}}{k!}$ (use the Ratio test) \vfill\newpage
\end{enumerate}
\textbf{HW problems 21-24. }\emph{Find the power series representations
centered at $0$ for the following functions using known power series.
Give the interval of convergence for the resulting series. }
\begin{enumerate}[resume]
\item $f(x)=\frac{1}{1-2x}$\vfill
\item $f(x)=\frac{2x^{3}}{3-3x}$\vfill
\end{enumerate}

\section*{11.1: Working with Parametric Equations.}

\textbf{HW problems 7-8. }\emph{Plot the parametric points for the
given values of $t$, and eliminate the parameter to an equation in
terms of $x$ and $y$. }
\begin{enumerate}[resume]
\item $x=-t+6$, $y=3t-3$; for $t=-5,-1,0,1,5$
\begin{figure}[H]
\begin{centering}
\includegraphics[width=2in]{slides_Calc3/images/empty_graph_-10-10_-20-20.png}
\par\end{centering}
\caption{Use the graph to complete the preceding problem.}

\end{figure}
\end{enumerate}
\textbf{HW problems 23-24 and 35-36. }\emph{Give a set of parametric
equations that describes the curves.}
\begin{enumerate}[resume]
\item The complete curve $x=y^{2}-2y$.\vfill
\item A man jogs clockwise at a constant speed around a circular track of
radius 50 yards completing one lap in 2.5 minutes.\vfill\newpage
\end{enumerate}
\textbf{HW problems 45-48. }\emph{Determine $\nicefrac{dy}{dx}$ in
terms of $t$ and evaluate it at the given value of $t$. }
\begin{enumerate}[resume]
\item $x=2t$, $y=t^{3}$; $t=-1$\vfill
\end{enumerate}

\section*{11.2: Working with Polar Coordinates.}

\textbf{HW problems 9-11. }\emph{Graph the points in the following
polar coordinates. }
\begin{enumerate}[resume]
\item $\left(1.5,\frac{7\pi}{4}\right)$, $\left(-1.5,\frac{7\pi}{4}\right)$,
$\left(1.5,-\frac{7\pi}{4}\right)$, $\left(1.5,\frac{7\pi}{4}+2\pi\right)$,
$\left(-1.5,\frac{7\pi}{4}+\pi\right)$ 
\begin{figure}[H]
\begin{centering}
\includegraphics[width=2in]{slides_Calc3/images/polar_graph_-2-2a.png}
\par\end{centering}
\caption{Use this graph to plot the polar coordinates}

\end{figure}
\end{enumerate}
\textbf{HW problems 15-17. }\emph{Express the following polar coordinates
in Cartesian coordinates.}
\begin{enumerate}[resume]
\item $\left(2,\frac{7\pi}{4}\right)$\vfill
\end{enumerate}
\textbf{HW problems 21-24. }\emph{Express the following Cartesian
coordinates in polar coordinates in at least two different ways. }
\begin{enumerate}[resume]
\item $\left(4,\sqrt{3}\right)$\vfill\newpage
\end{enumerate}
\textbf{HW problems 45. }\emph{Mark the points on the polar graph
that corresponds to the points shown on the Cartesian graph.}
\begin{enumerate}[resume]
\item 
\begin{figure}[H]
\begin{centering}
\emph{\includegraphics[width=2in]{slides_Calc3/images/graph_cos(t)+sin(2t).png}\includegraphics[width=2in]{slides_Calc3/images/polar_graph_cos(t)+sin(2t).png}}
\par\end{centering}
\caption{$r=\cos(\theta)+\sin(2\theta)$}

\end{figure}
Mark the following points: A, B, C, D, E, G, H, I, J
\end{enumerate}

\section*{11.3: Finding Derivatives and Integral in Polar Coordinates.}

\textbf{HW problems 5-11. }\emph{Find the slope of the line tangent
to the following polar curves at the given points. }
\begin{enumerate}[resume]
\item $1+\cos\theta$; $\left(1+\frac{\sqrt{3}}{2},\frac{\pi}{6}\right)$,
$\left(1,2\pi\right)$\vfill
\end{enumerate}
\textbf{HW problems 21-25 (odd). }\emph{Find the area of the region. }
\begin{enumerate}[resume]
\item The region inside the limacon $r=2+\cos\theta$.\vfill
\item The region inside one leaf of the rose $r=\cos5\theta$.\vfill
\end{enumerate}

\end{document}
